% page 399 of the facsimile (image p-398 of the scan)
% block 165.1 x 254.9 mm ; left margin 36.9 mm
\pgc{15.240mm}{0.000mm}{
\l{}
\l{}
\l{}
\l{}
\l{}
\l{}
\l{ne.\cel{2}Adverbo di nego.}
\l{}
\l{\sou{nebular}. (\sou{netrans}.) Dicesas kande aquo-vaporo nekomplete konden-}
\l{\cel{3}sita flotacas en ta parto di atmosfero qua esas apud tero}
\l{\cel{3}e trublas olua diafaneso. -\cel{2}DEFIS.}
\l{}
\l{\sou{nebulos}o. (\sou{astron.)}\cel{3}Aglomerajo ek steli qui esas tante multa}
\l{\cel{3}e semble inter-apuda ke li ne plus esas dicernebla singla.}
\l{\cel{3}-\cel{2}DEFIRS.}
\l{}
\l{\sou{necesa}.\cel{2}I.Quan onu ne povas karear. - II. De qua onu ne povas}
\l{\cel{3}dispensar su. - III. (\sou{filoz.}) Qua ne povas ne esar. -\cel{2}EFIS.}
\l{}
\l{\sou{negar}. (\sou{trans.})\cel{2}Deklarar ke ulo ne esas.\cel{2}DEFIS.}
\l{}
\l{\sou{negativa}.\cel{2}I. (\sou{algebro ed elektr.})\cel{2}Qua havas la signo minus. -}
\l{\cel{3}II. Negativo (\sou{fotogr.}) : Ube la nigro di la modelo aparas}
\l{\cel{3}blanka ; e la blanko, nigra. - DEFIRS.}
\l{}
\l{\sou{neglijar}\cel{1}(\sou{trans.)}\cel{3}I. Lasar (ulu, ulo) sen sorgar lu.}
\l{\cel{3}- II. Lasar sen egardar lu. - EFIS.}
\l{}
\l{\sou{neglijeo}.\cel{2}Homo-vesto,quan onu +weras (generale) dum la matino.}
\l{\cel{3}DEFI.}
\l{}
\l{\sou{negociar}. (\sou{trans.)}\cel{3}Mediacar por konkluzar afero (mariajo,}
\l{\cel{3}kontrato, paco). - EFIRS.}
\l{}
\l{\sou{negro}. Homo ek la raso pelnigra. DEFIS.}
\l{}
\l{\sou{nek}. Konjunciono qua separas propozicioni negala, o termini}
\l{\cel{3}interdiferanta, di propoziciono negala. FIRS.}
\l{}
\l{\sou{nekrobioso.}\cel{2}(\sou{patol.)}\cel{2}Modifikeso strukturala di organo o tisuo}
\l{\cel{3}en qua la sango-cirkulo cesis, generale pro infarkto. - DEFIS}
\l{}
\l{\sou{nekrologo.}\cel{3}Notico pri un o plura personegi jus mortinta.DEFIRS}
\l{}
\l{\sou{nekromanciar}\cel{2}(\sou{netrans.)}\cel{2}Praktikar la arto (asertita) konjurar}
\l{\cel{3}la mortinti per magio por obtenar de li la revelo di la event-}
\l{\cel{3}ontaji, di la kozi okulta.\cel{2}DEFIRS.}
\l{}
\l{nekropolo. I. (\sou{epoki antiqua)}\cel{2}Sorto di tombeyo subtera che}
\l{\cel{3}la Egiptiani ed altra populi antiqua. - II. Parto di urbo, qua}
\l{\cel{3}esas destinita a la sepulturi. - DEFIS.}
\l{}
\l{nekroso. - (\sou{patol}.) Stando di osto, o di parto de osto, qua}
\l{\cel{3}(che vivanto) cesis vivar. - DEFIS.}
\l{}
\l{nektaro.\cel{2}I. (\sou{mitol.}) Drinkajo delikatega di la dei. - II.}
\l{\cel{3}Drinkajo bonega, delikatega, qua ecelas irga ceteri. - III.}
\l{\cel{3}Liquido sukroza, quan sekrecas la folio o la floro, e qua}
\l{\cel{3}atraktas la insekti. - DEFIRS.}
}
